% GUID v6, Primer P8 (label box:stats): the statistics behind the scoring rules of Section sec:scoring. Background only; it states
% no result of the paper and does not repeat the rules themselves. Place it in Section sec:scoring: % GUID v6, Primer P8 (label box:stats): the statistics behind the scoring rules of Section sec:scoring. Background only; it states
% no result of the paper and does not repeat the rules themselves. Place it in Section sec:scoring: % GUID v6, Primer P8 (label box:stats): the statistics behind the scoring rules of Section sec:scoring. Background only; it states
% no result of the paper and does not repeat the rules themselves. Place it in Section sec:scoring: % GUID v6, Primer P8 (label box:stats): the statistics behind the scoring rules of Section sec:scoring. Background only; it states
% no result of the paper and does not repeat the rules themselves. Place it in Section sec:scoring: \input{primers/p8-stats}
\begin{primer}{Statistics}{box:stats}
\emph{Pulls and 95\%.} A measurement $x\pm\sigma$ usually means a Gaussian error with standard deviation $\sigma$. If a prediction
$x_G\pm\sigma_G$ is right and the two errors are independent, the difference $x-x_G$ is Gaussian with variance $\sigma^2+\sigma_G^2$, so the
\emph{pull} $p=(x-x_G)/\sqrt{\sigma^2+\sigma_G^2}$ is a standard normal number. It falls outside $\pm1.96$ with probability 5\% (two-sided),
so $|p|<1.96$ is a test at 95\%, and a published 95\% half-width $b$ corresponds to $\sigma=b/1.96$. \emph{Worked pull:} a model predicts
$10.0\pm0.3$ and an experiment finds $11.0\pm0.4$. Then $p=1.0/\sqrt{0.4^2+0.3^2}=1.0/0.5=2.0$: a FAIL at 95\%, although a right model
gives $|p|\ge2.0$ in 4.6\% of independent tests. For $n$ independent tests, about $0.05\,n$ FAILs are expected by chance, and the sum of the
squared pulls follows a $\chi^2$ distribution with $n$ degrees of freedom.

\emph{BIC and $\Delta$BIC.} A model with $k$ fitted parameters that reaches the maximum likelihood $\mathcal L_{\max}$ on $N$ data points has
$\mathrm{BIC}=-2\ln\mathcal L_{\max}+k\ln N$~\cite{schwarz1978}. The penalty $k\ln N$ charges each extra parameter. For large $N$ and
independent points, the difference $\Delta$BIC of two models approximates $-2$ times the logarithm of their Bayes factor, the ratio of the
probabilities of the data under the two models. A $\Delta$BIC of 2 to 6 is then positive evidence, 6 to 10 strong and above 10 very
strong~\cite{kassraftery1995}. The approximation assumes $N$ independent points. If the points come in groups that share an error, such
as the distance of one galaxy, the effective number of independent data is closer to the number of groups. A point-wise $\Delta$BIC then
overstates the evidence, roughly in proportion to the number of points per group.

\emph{The bootstrap.} To learn how much a fitted value would vary in a new sample, one draws many new samples of the same size from the
data themselves, with replacement, and refits each~\cite{efron1979}. The spread of the refitted values estimates the uncertainty; the 16th
to 84th percentiles bracket a range comparable to $\pm1\sigma$. The units resampled must be the independent ones, for rotation curves whole
galaxies, not single points.

\emph{Many tests.} With $m$ tests at 5\% each, the chance of at least one false positive, the \emph{family-wise error rate}, grows towards
$1-0.95^m$. Bonferroni's remedy tests each at $0.05/m$. Holm's method orders the $p$-values, $p_{(1)}\le\dots\le p_{(m)}$, and rejects
$p_{(i)}$ while $p_{(i)}\le0.05/(m-i+1)$; it keeps the family-wise error rate at 5\% under any dependence and never rejects less than
Bonferroni~\cite{holm1979}. Fisher's method combines $n$ \emph{independent} $p$-values: under the null each $p_i$ is uniform, so
$-2\ln p_i$ follows $\chi^2$ with 2 degrees of freedom, and $X^2=-2\sum_i\ln p_i$ follows $\chi^2$ with $2n$~\cite{fisher1932}. It responds
strongly to a few very small $p_i$.

\emph{Look-elsewhere and preregistration.} A search over many masses, scales or candidate values finds a small \emph{local} $p$-value
somewhere by chance. The \emph{global} $p$-value, the chance of so small a value anywhere in the search, is larger, roughly the local one
times a trials factor~\cite{grossvitells2010}. A test is confirmatory only if its hypothesis, data, statistic and decision rule were
fixed before the outcome was seen; \emph{preregistration} writes them down in a dated record beforehand, and every choice made afterwards
is then visible as exploratory~\cite{nosek2018}.
\end{primer}

\begin{primer}{Statistics}{box:stats}
\emph{Pulls and 95\%.} A measurement $x\pm\sigma$ usually means a Gaussian error with standard deviation $\sigma$. If a prediction
$x_G\pm\sigma_G$ is right and the two errors are independent, the difference $x-x_G$ is Gaussian with variance $\sigma^2+\sigma_G^2$, so the
\emph{pull} $p=(x-x_G)/\sqrt{\sigma^2+\sigma_G^2}$ is a standard normal number. It falls outside $\pm1.96$ with probability 5\% (two-sided),
so $|p|<1.96$ is a test at 95\%, and a published 95\% half-width $b$ corresponds to $\sigma=b/1.96$. \emph{Worked pull:} a model predicts
$10.0\pm0.3$ and an experiment finds $11.0\pm0.4$. Then $p=1.0/\sqrt{0.4^2+0.3^2}=1.0/0.5=2.0$: a FAIL at 95\%, although a right model
gives $|p|\ge2.0$ in 4.6\% of independent tests. For $n$ independent tests, about $0.05\,n$ FAILs are expected by chance, and the sum of the
squared pulls follows a $\chi^2$ distribution with $n$ degrees of freedom.

\emph{BIC and $\Delta$BIC.} A model with $k$ fitted parameters that reaches the maximum likelihood $\mathcal L_{\max}$ on $N$ data points has
$\mathrm{BIC}=-2\ln\mathcal L_{\max}+k\ln N$~\cite{schwarz1978}. The penalty $k\ln N$ charges each extra parameter. For large $N$ and
independent points, the difference $\Delta$BIC of two models approximates $-2$ times the logarithm of their Bayes factor, the ratio of the
probabilities of the data under the two models. A $\Delta$BIC of 2 to 6 is then positive evidence, 6 to 10 strong and above 10 very
strong~\cite{kassraftery1995}. The approximation assumes $N$ independent points. If the points come in groups that share an error, such
as the distance of one galaxy, the effective number of independent data is closer to the number of groups. A point-wise $\Delta$BIC then
overstates the evidence, roughly in proportion to the number of points per group.

\emph{The bootstrap.} To learn how much a fitted value would vary in a new sample, one draws many new samples of the same size from the
data themselves, with replacement, and refits each~\cite{efron1979}. The spread of the refitted values estimates the uncertainty; the 16th
to 84th percentiles bracket a range comparable to $\pm1\sigma$. The units resampled must be the independent ones, for rotation curves whole
galaxies, not single points.

\emph{Many tests.} With $m$ tests at 5\% each, the chance of at least one false positive, the \emph{family-wise error rate}, grows towards
$1-0.95^m$. Bonferroni's remedy tests each at $0.05/m$. Holm's method orders the $p$-values, $p_{(1)}\le\dots\le p_{(m)}$, and rejects
$p_{(i)}$ while $p_{(i)}\le0.05/(m-i+1)$; it keeps the family-wise error rate at 5\% under any dependence and never rejects less than
Bonferroni~\cite{holm1979}. Fisher's method combines $n$ \emph{independent} $p$-values: under the null each $p_i$ is uniform, so
$-2\ln p_i$ follows $\chi^2$ with 2 degrees of freedom, and $X^2=-2\sum_i\ln p_i$ follows $\chi^2$ with $2n$~\cite{fisher1932}. It responds
strongly to a few very small $p_i$.

\emph{Look-elsewhere and preregistration.} A search over many masses, scales or candidate values finds a small \emph{local} $p$-value
somewhere by chance. The \emph{global} $p$-value, the chance of so small a value anywhere in the search, is larger, roughly the local one
times a trials factor~\cite{grossvitells2010}. A test is confirmatory only if its hypothesis, data, statistic and decision rule were
fixed before the outcome was seen; \emph{preregistration} writes them down in a dated record beforehand, and every choice made afterwards
is then visible as exploratory~\cite{nosek2018}.
\end{primer}

\begin{primer}{Statistics}{box:stats}
\emph{Pulls and 95\%.} A measurement $x\pm\sigma$ usually means a Gaussian error with standard deviation $\sigma$. If a prediction
$x_G\pm\sigma_G$ is right and the two errors are independent, the difference $x-x_G$ is Gaussian with variance $\sigma^2+\sigma_G^2$, so the
\emph{pull} $p=(x-x_G)/\sqrt{\sigma^2+\sigma_G^2}$ is a standard normal number. It falls outside $\pm1.96$ with probability 5\% (two-sided),
so $|p|<1.96$ is a test at 95\%, and a published 95\% half-width $b$ corresponds to $\sigma=b/1.96$. \emph{Worked pull:} a model predicts
$10.0\pm0.3$ and an experiment finds $11.0\pm0.4$. Then $p=1.0/\sqrt{0.4^2+0.3^2}=1.0/0.5=2.0$: a FAIL at 95\%, although a right model
gives $|p|\ge2.0$ in 4.6\% of independent tests. For $n$ independent tests, about $0.05\,n$ FAILs are expected by chance, and the sum of the
squared pulls follows a $\chi^2$ distribution with $n$ degrees of freedom.

\emph{BIC and $\Delta$BIC.} A model with $k$ fitted parameters that reaches the maximum likelihood $\mathcal L_{\max}$ on $N$ data points has
$\mathrm{BIC}=-2\ln\mathcal L_{\max}+k\ln N$~\cite{schwarz1978}. The penalty $k\ln N$ charges each extra parameter. For large $N$ and
independent points, the difference $\Delta$BIC of two models approximates $-2$ times the logarithm of their Bayes factor, the ratio of the
probabilities of the data under the two models. A $\Delta$BIC of 2 to 6 is then positive evidence, 6 to 10 strong and above 10 very
strong~\cite{kassraftery1995}. The approximation assumes $N$ independent points. If the points come in groups that share an error, such
as the distance of one galaxy, the effective number of independent data is closer to the number of groups. A point-wise $\Delta$BIC then
overstates the evidence, roughly in proportion to the number of points per group.

\emph{The bootstrap.} To learn how much a fitted value would vary in a new sample, one draws many new samples of the same size from the
data themselves, with replacement, and refits each~\cite{efron1979}. The spread of the refitted values estimates the uncertainty; the 16th
to 84th percentiles bracket a range comparable to $\pm1\sigma$. The units resampled must be the independent ones, for rotation curves whole
galaxies, not single points.

\emph{Many tests.} With $m$ tests at 5\% each, the chance of at least one false positive, the \emph{family-wise error rate}, grows towards
$1-0.95^m$. Bonferroni's remedy tests each at $0.05/m$. Holm's method orders the $p$-values, $p_{(1)}\le\dots\le p_{(m)}$, and rejects
$p_{(i)}$ while $p_{(i)}\le0.05/(m-i+1)$; it keeps the family-wise error rate at 5\% under any dependence and never rejects less than
Bonferroni~\cite{holm1979}. Fisher's method combines $n$ \emph{independent} $p$-values: under the null each $p_i$ is uniform, so
$-2\ln p_i$ follows $\chi^2$ with 2 degrees of freedom, and $X^2=-2\sum_i\ln p_i$ follows $\chi^2$ with $2n$~\cite{fisher1932}. It responds
strongly to a few very small $p_i$.

\emph{Look-elsewhere and preregistration.} A search over many masses, scales or candidate values finds a small \emph{local} $p$-value
somewhere by chance. The \emph{global} $p$-value, the chance of so small a value anywhere in the search, is larger, roughly the local one
times a trials factor~\cite{grossvitells2010}. A test is confirmatory only if its hypothesis, data, statistic and decision rule were
fixed before the outcome was seen; \emph{preregistration} writes them down in a dated record beforehand, and every choice made afterwards
is then visible as exploratory~\cite{nosek2018}.
\end{primer}

\begin{primer}{Statistics}{box:stats}
\emph{Pulls and 95\%.} A measurement $x\pm\sigma$ usually means a Gaussian error with standard deviation $\sigma$. If a prediction
$x_G\pm\sigma_G$ is right and the two errors are independent, the difference $x-x_G$ is Gaussian with variance $\sigma^2+\sigma_G^2$, so the
\emph{pull} $p=(x-x_G)/\sqrt{\sigma^2+\sigma_G^2}$ is a standard normal number. It falls outside $\pm1.96$ with probability 5\% (two-sided),
so $|p|<1.96$ is a test at 95\%, and a published 95\% half-width $b$ corresponds to $\sigma=b/1.96$. \emph{Worked pull:} a model predicts
$10.0\pm0.3$ and an experiment finds $11.0\pm0.4$. Then $p=1.0/\sqrt{0.4^2+0.3^2}=1.0/0.5=2.0$: a FAIL at 95\%, although a right model
gives $|p|\ge2.0$ in 4.6\% of independent tests. For $n$ independent tests, about $0.05\,n$ FAILs are expected by chance, and the sum of the
squared pulls follows a $\chi^2$ distribution with $n$ degrees of freedom.

\emph{BIC and $\Delta$BIC.} A model with $k$ fitted parameters that reaches the maximum likelihood $\mathcal L_{\max}$ on $N$ data points has
$\mathrm{BIC}=-2\ln\mathcal L_{\max}+k\ln N$~\cite{schwarz1978}. The penalty $k\ln N$ charges each extra parameter. For large $N$ and
independent points, the difference $\Delta$BIC of two models approximates $-2$ times the logarithm of their Bayes factor, the ratio of the
probabilities of the data under the two models. A $\Delta$BIC of 2 to 6 is then positive evidence, 6 to 10 strong and above 10 very
strong~\cite{kassraftery1995}. The approximation assumes $N$ independent points. If the points come in groups that share an error, such
as the distance of one galaxy, the effective number of independent data is closer to the number of groups. A point-wise $\Delta$BIC then
overstates the evidence, roughly in proportion to the number of points per group.

\emph{The bootstrap.} To learn how much a fitted value would vary in a new sample, one draws many new samples of the same size from the
data themselves, with replacement, and refits each~\cite{efron1979}. The spread of the refitted values estimates the uncertainty; the 16th
to 84th percentiles bracket a range comparable to $\pm1\sigma$. The units resampled must be the independent ones, for rotation curves whole
galaxies, not single points.

\emph{Many tests.} With $m$ tests at 5\% each, the chance of at least one false positive, the \emph{family-wise error rate}, grows towards
$1-0.95^m$. Bonferroni's remedy tests each at $0.05/m$. Holm's method orders the $p$-values, $p_{(1)}\le\dots\le p_{(m)}$, and rejects
$p_{(i)}$ while $p_{(i)}\le0.05/(m-i+1)$; it keeps the family-wise error rate at 5\% under any dependence and never rejects less than
Bonferroni~\cite{holm1979}. Fisher's method combines $n$ \emph{independent} $p$-values: under the null each $p_i$ is uniform, so
$-2\ln p_i$ follows $\chi^2$ with 2 degrees of freedom, and $X^2=-2\sum_i\ln p_i$ follows $\chi^2$ with $2n$~\cite{fisher1932}. It responds
strongly to a few very small $p_i$.

\emph{Look-elsewhere and preregistration.} A search over many masses, scales or candidate values finds a small \emph{local} $p$-value
somewhere by chance. The \emph{global} $p$-value, the chance of so small a value anywhere in the search, is larger, roughly the local one
times a trials factor~\cite{grossvitells2010}. A test is confirmatory only if its hypothesis, data, statistic and decision rule were
fixed before the outcome was seen; \emph{preregistration} writes them down in a dated record beforehand, and every choice made afterwards
is then visible as exploratory~\cite{nosek2018}.
\end{primer}
